\documentclass[10pt,twocolumn]{article}
\usepackage[T1]{fontenc}
\usepackage[utf8]{inputenc}
\usepackage{lmodern}
\usepackage[margin=1.8cm,top=2.4cm,bottom=2cm,columnsep=0.6cm]{geometry}
\usepackage{fancyhdr}
\usepackage{titlesec}
\usepackage{booktabs}
\usepackage{amsmath,amssymb,amsfonts}
\usepackage{microtype}
\usepackage{xcolor}
\usepackage{mdframed}
\usepackage{parskip}
\usepackage{enumitem}
\usepackage{hyperref}

\definecolor{rulered}{RGB}{180,30,30}
\definecolor{boxgray}{RGB}{245,245,245}
\definecolor{keywordgray}{RGB}{100,100,100}

\pagestyle{fancy}
\fancyhf{}
\fancyhead[L]{\textsc{\small Claude Mag}}
\fancyhead[C]{\small\textit{Machine Learning Research Digest}}
\fancyhead[R]{\small 2026-11-09}
\fancyfoot[C]{\small\thepage}
\renewcommand{\headrulewidth}{0.8pt}

\titleformat{\section}{\normalsize\bfseries\color{rulered}}{}{0pt}{}%
  [\vspace{-4pt}\color{rulered}\rule{\columnwidth}{0.4pt}]
\titlespacing{\section}{0pt}{8pt}{4pt}

\hypersetup{colorlinks=true,urlcolor=rulered,linkcolor=black}

\begin{document}
\setlength{\parindent}{0pt}
\setlength{\parskip}{4pt}

{\small\color{keywordgray}\textbf{Keywords:}
  evolutionary optimization \textbullet{}
  latent agents \textbullet{}
  program synthesis \textbullet{}
  scientific computing}\par\vspace{4pt}

{\LARGE\bfseries\raggedright
  SwarmEvolve: Latent Agent Swarms for Evolutionary Discovery
  and Optimization}\par\vspace{4pt}

{\large\itshape\raggedright
  Agents That Reason in Activation Space Rather Than Text
  Evaluate 5--13$\times$ More Programs Per Hour on Mathematical
  and Scientific Discovery}\par\vspace{6pt}

{\small\textbf{By} Qile Jiang, George Em Karniadakis}\par\vspace{2pt}

{\small\color{keywordgray}arXiv:2611.08483 \textbullet{}
  cs.AI / cs.NE \textbullet{} Brown University}
\par\vspace{2pt}\noindent{\color{rulered}\rule{\columnwidth}{0.6pt}}\vspace{6pt}

LLM-guided evolutionary search is accurate but slow because each
candidate requires a full chain-of-thought forward pass. SwarmEvolve
replaces text-decoding agents with latent agents that share KV caches
and reason in activation space, evaluating 5--13$\times$ more programs
per hour than state-of-the-art frameworks and reaching equivalent
solution quality in 3--8 minutes versus 60 minutes for competing systems.

\begin{mdframed}[backgroundcolor=boxgray,linecolor=rulered,linewidth=1pt,
    innerleftmargin=6pt,innerrightmargin=6pt,
    innertopmargin=6pt,innerbottommargin=6pt]
\textbf{\small AT A GLANCE}\par\vspace{4pt}
\begin{description}[leftmargin=0pt,labelsep=4pt,itemsep=2pt,parsep=0pt,
    font=\small\bfseries]
  \item[Problem:] LLM-guided evolutionary search is accurate but
    too slow for large-scale discovery; each round requires full
    chain-of-thought for every candidate.
  \item[Why it matters:] Scientific computing and mathematical
    discovery require evaluating hundreds of program variants per
    experiment; throughput is a hard bottleneck.
  \item[Method:] 10 swarms of 45 latent agents share KV caches and
    reason for 40 latent steps without decoding to text; a lightweight
    decoding head maps top agent states to program tokens.
  \item[Result:] 5--13$\times$ more programs per hour; solution quality
    reached in 3--8 min vs.\ 60 min for text-based frameworks.
  \item[Evidence:] 21 problems spanning symbolic math and PDE solver
    kernel optimization; baselines: OpenEvolve, ShinkaEvolve, CodeEvolve.
\end{description}
\end{mdframed}

\section{The Big Picture}

Program evolution with LLMs (OpenEvolve, FunSearch, ShinkaEvolve) has
produced surprising discoveries in combinatorics and numerical computing
by using LLMs as intelligent mutation operators: given a population of
programs and their fitness scores, the LLM proposes new variants that
improve on the best candidates. The approach works because LLMs encode
prior knowledge about program structure, algorithm design, and
mathematical transformations that random mutation cannot access.

The fundamental bottleneck is throughput. Generating a new program
candidate with chain-of-thought reasoning costs one full LLM forward
pass with a long context window. At 450 tokens per second on typical
hardware, a 1,000-token chain-of-thought allows only 27 candidates per
minute per GPU. Scientific computing problems may need thousands of
candidates before a significant improvement is found.

SwarmEvolve attacks this bottleneck by eliminating the decoding step
from most of the computation. Agents in SwarmEvolve do not produce
natural language --- they produce latent representations and hand them
to a compact decoding head. The decoding head is cheap, while the
latent reasoning is parallelized across 450 agents per round.

\section{Why Existing Approaches Fall Short}

\textbf{Text-based evolutionary LLM frameworks} (OpenEvolve, CodeEvolve)
generate each candidate as a full text completion. This is flexible but
slow: each candidate occupies the full generation path of the LLM.

\textbf{Parallel sampling with temperature} can explore multiple
candidates simultaneously, but each sample is independent and cannot
condition on what other concurrent samples are producing, limiting
diversity within a batch.

\textbf{Speculative decoding} speeds up text generation but does not
reduce the number of distinct candidates that can be evaluated per unit
time, since it optimizes throughput on a single sequence.

\textbf{Smaller surrogate models} are faster but lack the reasoning
quality of large LLMs, leading to lower-quality proposals and slower
convergence in solution quality.

\section{Core Mechanism: Latent Agents with KV Sharing}

Each agent is a frozen LLM conditioned on a shared prefix KV cache
covering the problem description and global best solution. The agent
then runs 40 autoregressive steps \emph{without} applying the
unembedding matrix: no token is produced, and no text is decoded.
Instead, the final hidden state $\mathbf{h}_i \in \mathbb{R}^d$
encodes the agent's ``proposal'' in activation space.

Agents in the same swarm share a common prefix KV cache $\mathcal{K}$,
which eliminates redundant computation for the shared context and
implicitly exposes each agent to the best solution found so far.
Agents have different random seeds (different initial noise vectors
appended to the shared prefix), inducing diversity within the swarm.

A lightweight decoding head $D_\phi: \mathbb{R}^d \to \mathcal{V}^L$
maps the hidden states of the top-$r$ agents (by cosine similarity to
the fitness-weighted centroid of the swarm) to program token sequences
of length at most $L = 200$ tokens.

\section{Technical Formulation}

Let $\mathcal{S} = \{S_1, \ldots, S_K\}$ denote $K = 10$ swarms, each
containing $N = 45$ agents. The shared prefix for swarm $S_k$ is:
\[
  \mathcal{K}_k = \text{KVCache}\bigl([\text{problem}, \,
    \mathbf{x}^*_k]\bigr),
\]
where $\mathbf{x}^*_k$ is the globally best program found in rounds
$1, \ldots, t-1$ of swarm $k$.

Each agent $a_{k,i}$ produces a hidden state:
\[
  \mathbf{h}_{k,i} = f_\theta\bigl(\mathbf{z}_{k,i},
    \mathcal{K}_k\bigr) \in \mathbb{R}^d,
\]
where $\mathbf{z}_{k,i} \sim \mathcal{N}(0, \sigma^2 I)$ is the
agent-specific noise vector and $f_\theta$ runs 40 latent steps.

The swarm centroid weighted by fitness is:
\[
  \bar{\mathbf{h}}_k = \frac{\sum_{i=1}^N w_i \mathbf{h}_{k,i}}
    {\sum_{i=1}^N w_i}, \quad
  w_i = \text{softmax}(F_{k,i} / T),
\]
where $F_{k,i}$ is the previous-round fitness of agent $a_{k,i}$ and
$T$ is a temperature parameter. The top-$r = 5$ agents (by cosine
similarity to $\bar{\mathbf{h}}_k$) are decoded by $D_\phi$.

\section{Learning or Inference Procedure}

$D_\phi$ is a 2-layer MLP trained offline on (hidden-state, program)
pairs extracted from text-based evolutionary runs. At inference time,
$D_\phi$ is frozen; SwarmEvolve requires no online training.

The evolutionary loop per round:
\begin{enumerate}[itemsep=2pt,parsep=0pt]
  \item Update all $K$ prefix caches with current best solutions.
  \item Run all $K \times N = 450$ agents in parallel (each 40 latent
    steps) on a single A100 node.
  \item Select top-$r$ agents per swarm using the fitness-weighted
    cosine similarity score.
  \item Decode selected agents' hidden states to programs via $D_\phi$.
  \item Execute decoded programs in a sandboxed environment.
  \item Update global best solutions for each swarm.
\end{enumerate}

\section{What the Guarantee Says}

No formal convergence guarantee is provided. The empirical claim is
throughput superiority on 21 specific problems. The best-known
solutions to 11 of the 21 problems (symbolic math subset) have formal
certificates; SwarmEvolve recovers certified-correct solutions for 9 of
these 11 within the allocated time budget of 2 GPU-hours.

\section{Experimental Findings}

\begin{itemize}[itemsep=2pt,parsep=0pt]
  \item \textbf{Throughput:} 5--13$\times$ more programs per hour than
    all three baselines on all 21 problems.
  \item \textbf{Time-to-solution:} Median 3--8 minutes to reach the
    quality of the best competing solution; baselines require 60 minutes
    on the same hardware.
  \item \textbf{Solution quality:} Highest average score across 21
    problems; wins 17 of 21 head-to-head comparisons vs.\ the best
    single baseline per problem.
  \item \textbf{Diversity:} The 10-swarm configuration produces
    1.9$\times$ more distinct program families per hour than a
    450-agent single-swarm variant, confirming the value of swarm
    isolation for population diversity.
\end{itemize}

\section{Ablations and Interpretation}

\textbf{Latent steps:} 40 latent steps is near-optimal; 20 steps
reduces quality by 4.3 average points; 80 steps shows no improvement
while doubling per-round time.

\textbf{Swarm count:} 10 swarms outperform 5 swarms by 2.7 average
points and 1 swarm by 7.1 points, confirming that swarm diversity is
the main quality driver.

\textbf{KV sharing:} Removing the shared prefix KV cache (agents start
from scratch each round) reduces throughput by 2.3$\times$ and quality
by 3.2 points, since agents lose access to the current best solution.

\textbf{Decoding head quality:} Replacing $D_\phi$ with a linear probe
reduces program validity (fraction of programs that compile and run)
from 94\% to 71\%, confirming the need for a non-linear decoder.

\section*{Reference}

Qile Jiang, George Em Karniadakis.
\textbf{SwarmEvolve: Latent Agent Swarms for Evolutionary Discovery
and Optimization.}
arXiv:2611.08483.
\url{https://arxiv.org/abs/2611.08483}

\end{document}
